% Borrador público del TFG de Física (versión en español).
% Compilar: latexmk -pdf main.tex
\documentclass[11pt,a4paper,oneside]{report}
% Preámbulo de la versión en español.

\usepackage[utf8]{inputenc}
\usepackage[T1]{fontenc}
\usepackage[spanish,es-noshorthands,es-tabla]{babel}

\usepackage[a4paper,top=2.5cm,bottom=2.5cm,left=2.5cm,right=2.5cm]{geometry}
\usepackage{setspace}
\onehalfspacing

\usepackage{amsmath,amssymb,amsthm,mathtools}
\usepackage{bm}

\usepackage{graphicx}
\usepackage{xcolor}
\usepackage{tikz}
\usetikzlibrary{arrows.meta,positioning,calc}
\usepackage{enumitem}
\usepackage{booktabs}
\usepackage{caption}
\usepackage[nottoc]{tocbibind}
\usepackage[colorlinks=true,linkcolor=black,citecolor=blue,urlcolor=blue]{hyperref}
\hypersetup{
  pdftitle={Mecánica cuántica euclídea aplicada a la valoración de opciones path-dependent (borrador)},
  pdfauthor={Roberto Cepeda Rocandio}
}

\theoremstyle{plain}
\newtheorem{teorema}{Teorema}[chapter]
\newtheorem{proposicion}[teorema]{Proposición}
\newtheorem{lema}[teorema]{Lema}
\newtheorem{corolario}[teorema]{Corolario}

\theoremstyle{definition}
\newtheorem{definicion}[teorema]{Definición}
\newtheorem{ejemplo}[teorema]{Ejemplo}

\theoremstyle{remark}
\newtheorem{observacion}[teorema]{Observación}

\newcommand{\R}{\mathbb{R}}
\newcommand{\N}{\mathbb{N}}
\newcommand{\Z}{\mathbb{Z}}
\newcommand{\C}{\mathbb{C}}
\newcommand{\E}{\mathbb{E}}
\newcommand{\PP}{\mathbb{P}}
\newcommand{\Q}{\mathbb{Q}}
\newcommand{\F}{\mathcal{F}}
\newcommand{\G}{\mathcal{G}}
\newcommand{\Hh}{\mathcal{H}}
\newcommand{\HH}{\mathcal{H}}
\newcommand{\1}{\mathbf{1}}
\newcommand{\dd}{\mathrm{d}}
\newcommand{\eqdef}{\coloneqq}
\DeclareMathOperator{\Var}{Var}

\newcommand{\bra}[1]{\langle #1 |}
\newcommand{\ket}[1]{| #1 \rangle}
\newcommand{\braket}[2]{\langle #1 \mid #2 \rangle}
\newcommand{\matel}[3]{\langle #1 \mid #2 \mid #3 \rangle}
\newcommand{\Hbs}{H_{\text{BS}}}
\newcommand{\Heff}{H_{\text{eff}}}
\newcommand{\Lag}{\mathcal{L}}
\newcommand{\Dpath}{\mathcal{D}}

% Borrador en curso: remisiones y capítulos aún no publicados.
\newcommand{\enprep}{, en preparación}
\newif\ifprimerprevisto \primerprevistotrue
\newcommand{\capituloprevisto}[2]{%
  \ifprimerprevisto\clearpage\global\primerprevistofalse\else\par\vspace{2.5em}\fi
  \begingroup\let\clearpage\relax\let\cleardoublepage\relax
  \chapter{#1}\label{#2}\endgroup
  \noindent\textit{Capítulo en preparación.}\par}
\newcommand{\seccionesprevistas}[1]{%
  \par\medskip\noindent\textit{Secciones #1: en preparación.}\par}

\usepackage{titlesec}
\titleformat{\chapter}[hang]{\normalfont\huge\bfseries}{\thechapter.}{0.6em}{}
\titlespacing*{\chapter}{0pt}{-20pt}{18pt}


\newcommand{\versionborrador}{v0.1}
\newcommand{\fechaborrador}{17 de octubre de 2026}

\begin{document}

\pagenumbering{alph}
\begin{titlepage}
\centering
\vspace*{2cm}
{\LARGE\bfseries Mecánica cuántica euclídea aplicada\\[0.3cm]
a la valoración de opciones path-dependent\par}
\vspace{1.5cm}
{\Large Roberto Cepeda Rocandio\par}
\vspace{2.5cm}
{\large Trabajo Fin de Grado en Física\par}
\vspace{0.4cm}
{\sc Universidad de Oviedo}\\[0.2cm]
Facultad de Ciencias
\vfill
{\large\bfseries Borrador en curso\par}
\vspace{0.3cm}
Versión \versionborrador{} \textperiodcentered{} \fechaborrador\par
\vspace{0.6cm}
\begin{minipage}{0.8\textwidth}\centering\small
Los capítulos se publican como borradores a medida que se revisan.
Los que figuran como «en preparación» aún no se han publicado.
\end{minipage}
\end{titlepage}

\clearpage
\pagenumbering{roman}
\tableofcontents

\clearpage
\pagenumbering{arabic}
\setcounter{page}{1}

\chapter{Introducción}\label{cap:intro}

\section{Motivación}
\label{sec:motivacion}

En el modelo de Black--Scholes \cite{black-scholes, merton1973}, el
precio de una opción europea satisface una ecuación parabólica
retrógrada con el pago como condición final. Escrita en el logaritmo
del precio, $x = \log S$, y en el tiempo que falta hasta el vencimiento,
$\tau = T - t$, es una ecuación de Schrödinger en tiempo imaginario,
$\partial_\tau V = -\Hbs V$, para el hamiltoniano $\Hbs$ de una
partícula en una dimensión (Cap.~\ref{cap:bs_operadores}\enprep). El
precio es entonces un elemento de matriz del operador de evolución
euclídea $e^{-\tau\Hbs}$ integrado contra el pago, y la fórmula de
Feynman--Kac \cite{kac1949} lo escribe como una integral de camino
\cite{feynman1948} sobre las trayectorias del log-precio. La rotación de
Wick \cite{wick1954}, $t \to -i\tau$, que lleva de la mecánica cuántica
a la difusión, no hay que hacerla: el tiempo hasta el vencimiento es, de
origen, un tiempo euclídeo (Cap.~\ref{cap:fundamentos}\enprep).

\medskip

Para una opción europea, este lenguaje reproduce un resultado conocido.
Su interés está en las opciones cuyo pago depende de toda la trayectoria
del precio (\emph{path-dependent}). Su precio sigue siendo el valor
esperado del pago descontado, pero calcularlo exige la ley de una
funcional de la trayectoria, y cada familia de opciones pide una técnica
propia (Cap.~\ref{cap:problema_financiero}). En el lenguaje
hamiltoniano, la dependencia de la trayectoria entra en el mismo
operador de dos maneras: como un potencial o como una condición de
contorno. Ordenadas por la técnica que las resuelve, las opciones que
estudiará el trabajo recorren la progresión de un curso de mecánica
cuántica: partícula libre (europea), fuente lineal (asiática
geométrica), pared (barrera única), caja (doble barrera) y pared
exponencial (asiática aritmética, potencial de Liouville).

\medskip

El problema tiene además interés físico propio. En tiempo euclídeo, la
integral de camino está respaldada por una medida rigurosa, la de
Wiener, mientras que en tiempo real no define una medida; la rotación de
Wick es la continuación analítica que convierte el grupo unitario
$e^{-iHt}$ en el semigrupo $e^{-\tau H}$
(Cap.~\ref{cap:fundamentos}\enprep). El hamiltoniano de Black--Scholes
no es hermítico, y una transformación gauge, que resulta ser el teorema
de Girsanov escrito en operadores, lo conjuga con una partícula libre
más una constante (Cap.~\ref{cap:bs_operadores}\enprep). La asiática
aritmética conduce a la mecánica cuántica de Liouville, con espectro
continuo y funciones de Macdonald de índice imaginario como
autofunciones (Cap.~\ref{cap:espectro}\enprep). Y como en tiempo
euclídeo el peso de cada trayectoria es positivo, la integral de camino
discretizada se puede simular por Monte Carlo sin el problema del signo
del tiempo real (Cap.~\ref{cap:numerica}\enprep).

\section{Estado del arte}
\label{sec:estado_arte}

\noindent\textbf{Física y valoración de derivados.}
La aplicación de métodos de la física a los mercados tiene dos
vertientes. La física estadística ha estudiado sobre todo la estadística
empírica de los precios, con colas gruesas y correlaciones que el modelo
de Black--Scholes no recoge, y sus consecuencias para el riesgo y la
cobertura \cite{bouchaud-potters}. La otra conserva el modelo y cambia
el lenguaje. Linetsky \cite{linetsky1998} presentó a una audiencia
financiera la valoración como integral de camino; Baaquie
\cite{baaquie2004} la desarrolló de forma sistemática con los
hamiltonianos de la mecánica cuántica, y la extendió a la teoría de
campos para los tipos de interés \cite{baaquie2018}; Kleinert
\cite{kleinert} dedica a los mercados financieros una parte de su
tratado de integrales de camino. Este trabajo pertenece a la segunda
vertiente.

\medskip

\noindent\textbf{Resultados por familia.}
Todos los precios que se estudiarán son conocidos. La fórmula de la
asiática geométrica es de Kemna y Vorst \cite{kemna-vorst}; la de la
barrera única, de Merton \cite{merton1973}, y las de sus ocho variantes,
de Reiner y Rubinstein \cite{reiner-rubinstein}; la serie de imágenes de
la doble barrera, de Kunitomo e Ikeda \cite{kunitomo-ikeda}. La asiática
aritmética es el caso difícil: Yor \cite{yor1992} obtuvo la ley de la
integral de la exponencial del browniano; Geman y Yor \cite{geman-yor},
la transformada de Laplace del precio mediante procesos de Bessel, y
Linetsky \cite{linetsky2004}, su expansión espectral, que es la
formulación más próxima a la de este trabajo. En la parte numérica,
Glasserman \cite{glasserman} es la referencia de Monte Carlo en
finanzas, y Broadie, Glasserman y Kou \cite{broadie-glasserman-kou}
cuantificaron el efecto de vigilar la barrera solo en fechas discretas.

\medskip

\noindent\textbf{Base matemática.}
El cálculo estocástico está en Karatzas y Shreve \cite{karatzas-shreve}
y en Øksendal \cite{oksendal}, y su aplicación a la valoración, en
Shreve \cite{shreve}, Hull \cite{hull} y el TFG de Matemáticas
\cite{tfgm}; la integral de camino, en Feynman y Hibbs
\cite{feynman-hibbs}; la teoría espectral, en Reed y Simon
\cite{reed-simon1, reed-simon2} y Simon \cite{simon-functional}, y las
funciones especiales, en Lebedev \cite{lebedev}.

\medskip

La aportación que se busca no está, por tanto, en los resultados, sino
en su organización: obtenerlos todos como propagadores o espectros de un
mismo hamiltoniano, ordenados por la técnica física que los resuelve, y
contrastarlos numéricamente.

\section{Objetivos}
\label{sec:objetivos}

El trabajo tiene cuatro objetivos:
\begin{enumerate}[leftmargin=*, label=\arabic*., itemsep=1pt, topsep=3pt]
 \item \textbf{Desarrollar el formalismo con rigor:} difusión, medida de
 Wiener, integral de camino de Feynman, rotación de Wick, fórmula de
 Feynman--Kac y formalismo hamiltoniano, señalando en cada caso qué es
 riguroso y qué es formal.
 \item \textbf{Validar el marco} recuperando la fórmula de
 Black--Scholes como el propagador de una partícula libre.
 \item \textbf{Aplicarlo a opciones \emph{path-dependent}, ordenadas por
 técnica:} fuente lineal (asiática geométrica), pared absorbente
 (barrera única), caja (doble barrera) y potencial de Liouville
 (asiática aritmética).
 \item \textbf{Contrastar los resultados numéricamente,} con Monte Carlo
 de la integral de camino discretizada y con la evaluación de las
 representaciones espectrales.
\end{enumerate}

\noindent\textbf{Alcance.}
El trabajo se ciñe al modelo de Black--Scholes estándar: volatilidad y
tipo de interés constantes, sin saltos y con monitorización continua.
Las opciones \emph{lookback} no se desarrollarán porque no aportan una
técnica nueva: son barreras cuyo nivel se integra
(Cap.~\ref{cap:fuente_pared}\enprep). Quedan fuera la volatilidad
estocástica, los saltos, las opciones americanas y los modelos de tipos
de interés, posibles líneas de continuación
(Cap.~\ref{cap:conclusiones}\enprep).

\section{Estructura de la memoria}
\label{sec:estructura_memoria}

El Cap.~\ref{cap:problema_financiero} planteará el problema financiero
en lenguaje físico: los contratos, el no arbitraje, la probabilidad
neutral al riesgo, el modelo y la ecuación de Black--Scholes, y por qué
las opciones \emph{path-dependent} necesitan otro método. El
Cap.~\ref{cap:fundamentos} construirá el aparato (objetivo~1), y el
Cap.~\ref{cap:bs_operadores} lo aplicará a la opción europea, que tras
la transformación gauge es una partícula libre (objetivo~2). Los
Caps.~\ref{cap:fuente_pared} y~\ref{cap:espectro} abordarán el
objetivo~3 en orden de dificultad técnica: en el primero bastarán
completar el cuadrado y el método de imágenes; en el segundo habrá que
diagonalizar el hamiltoniano, de espectro discreto en la caja y continuo
en el potencial de Liouville. El Cap.~\ref{cap:numerica} (objetivo~4)
simulará por Monte Carlo la integral de camino discretizada y la
contrastará con las fórmulas exactas. Estudiará la pared en el retículo,
donde vigilar la barrera solo en los nodos introduce un error que
decrece como $\sqrt{\Delta t}$ y que admite dos correcciones distintas:
la probabilidad de cruce del puente browniano entre dos nodos, que hace
que la simulación valore la barrera continua, y el desplazamiento de la
barrera de Broadie, Glasserman y Kou, que adapta la fórmula continua a
una vigilancia discreta. Comparará además la serie espectral y la de
imágenes de la doble barrera según el vencimiento, y valorará la
asiática aritmética con la geométrica como variable de control. El
Cap.~\ref{cap:conclusiones} recogerá las conclusiones. Los
Caps.~\ref{cap:problema_financiero}--\ref{cap:bs_operadores} son
prerrequisito de los demás.

\section{Notación y convenios}
\label{sec:notacion_convenios}

La tabla siguiente recoge los símbolos principales.
\begin{center}
\footnotesize
\renewcommand{\arraystretch}{1.15}
\setlength{\tabcolsep}{4pt}
\begin{tabular}{@{}ll|ll@{}}
\toprule
$\PP,\ (\F_t)$ & Probabilidad real, filtración &
$\ket{\psi},\ \braket{\phi}{\psi}$ & Ket, producto escalar \\
$\Q,\ \E = \E_\Q$ & Medida neutral al riesgo &
$\ket{x},\ \ket{p}$ & Estados de posición y momento \\
$W_t,\ \widetilde W_t$ & Brownianos bajo $\PP$ y $\Q$ &
$H,\ \Hbs,\ \Heff$ & Hamiltonianos \\
$S_t,\ X_t = \log S_t$ & Precio y log-precio &
$\matel{x}{e^{-\tau H}}{x'}$ & Núcleo euclídeo \\
$r,\ \sigma,\ K,\ T$ & Tipo, volatilidad, \emph{strike}, vencimiento &
$\Dpath x,\ \Lag$ & Medida formal, lagrangiana \\
$\tau = T - t$ & Tiempo hasta el vencimiento &
$N(\cdot)$ & Normal estándar acumulada \\
$G_T,\ A_T$ & Promedios geométrico y aritmético &
$A_h^{(\nu)},\ \nu = \frac{2b}{\sigma^2}$ & Funcional de Yor e índice \\
$m_\tau$ & Mínimo del log-precio &
$B,\ \ell = \log B$ & Barrera ($\ell_d,\ \ell_u$) \\
$\mu,\ b = r - \frac{\sigma^2}{2}$ & Deriva real de $S$; de $X$ bajo $\Q$ &
$\theta;\ \beta = \frac{1}{2} - \frac{r}{\sigma^2}$
 & Parámetro gauge; su valor en BS \\
\bottomrule
\end{tabular}
\end{center}

\noindent Se usan unidades con $\hbar = 1$, de modo que el operador de
evolución euclídea es $e^{-\tau H}$. Toda esperanza sin subíndice es
bajo $\Q$. El parámetro gauge general $\theta$ toma en el modelo de
Black--Scholes el valor
$\beta = \frac{1}{2} - \frac{r}{\sigma^2} = -\frac{b}{\sigma^2}$. Se
sigue el convenio de Baaquie \cite{baaquie2004},
$\braket{x}{p} = e^{ipx}$, con las relaciones de completitud
$\int dx\,\ket{x}\bra{x} = \1 = \int \frac{dp}{2\pi}\,\ket{p}\bra{p}$.
Las demás constantes adimensionales y los reescalados temporales se
introducirán donde aparezcan.

\capituloprevisto{El problema financiero}{cap:problema_financiero}

\capituloprevisto{Difusión, integral de camino y tiempo euclídeo}{cap:fundamentos}

\capituloprevisto{El hamiltoniano de Black--Scholes: partícula libre}{cap:bs_operadores}

\capituloprevisto{Fuente lineal y pared absorbente}{cap:fuente_pared}

\capituloprevisto{Espectro: caja y potencial de Liouville}{cap:espectro}

\capituloprevisto{Comprobación numérica}{cap:numerica}

\capituloprevisto{Conclusiones y perspectivas}{cap:conclusiones}


\clearpage
\begin{thebibliography}{99}

\bibitem{black-scholes}
F. Black, M. Scholes.
The pricing of options and corporate liabilities.
\emph{Journal of Political Economy}, 81 (3), 637--654, 1973.

\bibitem{merton1973}
R. C. Merton.
Theory of rational option pricing.
\emph{Bell Journal of Economics and Management Science}, 4 (1),
141--183, 1973.

\bibitem{baaquie2004}
B. E. Baaquie.
\emph{Quantum Finance: Path Integrals and Hamiltonians for Options
and Interest Rates}.
Cambridge University Press, Cambridge, 2004.

\bibitem{baaquie2018}
B. E. Baaquie.
\emph{Quantum Field Theory for Economics and Finance}.
Cambridge University Press, Cambridge, 2018.

\bibitem{linetsky1998}
V. Linetsky.
The path integral approach to financial modeling and option pricing.
\emph{Computational Economics}, 11, 129--163, 1998.

\bibitem{linetsky2004}
V. Linetsky.
Spectral expansions for Asian (average price) options.
\emph{Operations Research}, 52 (6), 856--867, 2004.

\bibitem{kemna-vorst}
A. G. Z. Kemna, A. C. F. Vorst.
A pricing method for options based on average asset values.
\emph{Journal of Banking and Finance}, 14 (1), 113--129, 1990.

\bibitem{geman-yor}
H. Geman, M. Yor.
Bessel processes, Asian options, and perpetuities.
\emph{Mathematical Finance}, 3 (4), 349--375, 1993.

\bibitem{yor1992}
M. Yor.
On some exponential functionals of Brownian motion.
\emph{Advances in Applied Probability}, 24 (3), 509--531, 1992.

\bibitem{kunitomo-ikeda}
N. Kunitomo, M. Ikeda.
Pricing options with curved boundaries.
\emph{Mathematical Finance}, 2 (4), 275--298, 1992.

\bibitem{reiner-rubinstein}
E. Reiner, M. Rubinstein.
Breaking down the barriers.
\emph{Risk}, 4 (8), 28--35, 1991.

\bibitem{feynman1948}
R. P. Feynman.
Space-time approach to non-relativistic quantum mechanics.
\emph{Reviews of Modern Physics}, 20 (2), 367--387, 1948.

\bibitem{kac1949}
M. Kac.
On distributions of certain Wiener functionals.
\emph{Transactions of the American Mathematical Society}, 65 (1),
1--13, 1949.

\bibitem{wick1954}
G. C. Wick.
Properties of Bethe--Salpeter wave functions.
\emph{Physical Review}, 96 (4), 1124--1134, 1954.

\bibitem{kleinert}
H. Kleinert.
\emph{Path Integrals in Quantum Mechanics, Statistics, Polymer
Physics, and Financial Markets}.
World Scientific, Singapur, 5\textsuperscript{a} edición, 2009.

\bibitem{feynman-hibbs}
R. P. Feynman, A. R. Hibbs.
\emph{Quantum Mechanics and Path Integrals}.
McGraw-Hill, Nueva York, 1965.

\bibitem{reed-simon1}
M. Reed, B. Simon.
\emph{Methods of Modern Mathematical Physics I: Functional Analysis}.
Edición revisada, Academic Press, Nueva York, 1980.

\bibitem{reed-simon2}
M. Reed, B. Simon.
\emph{Methods of Modern Mathematical Physics II: Fourier Analysis,
Self-Adjointness}.
Academic Press, Nueva York, 1975.

\bibitem{simon-functional}
B. Simon.
\emph{Functional Integration and Quantum Physics}.
2\textsuperscript{a} edición, AMS Chelsea Publishing, Providence, 2005.

\bibitem{lebedev}
N. N. Lebedev.
\emph{Special Functions and Their Applications}.
Dover Publications, Nueva York, 1972.

\bibitem{karatzas-shreve}
I. Karatzas, S. E. Shreve.
\emph{Brownian Motion and Stochastic Calculus}.
2\textsuperscript{a} edición, Graduate Texts in Mathematics,
vol.~113. Springer, Nueva York, 1991.

\bibitem{shreve}
S. E. Shreve.
\emph{Stochastic Calculus for Finance II: Continuous-Time Models}.
Springer Finance. Springer, Nueva York, 2004.

\bibitem{oksendal}
B. Øksendal.
\emph{Stochastic Differential Equations: An Introduction with
Applications}.
6\textsuperscript{a} edición, Springer, Berlín, 2003.

\bibitem{hull}
J. C. Hull.
\emph{Options, Futures, and Other Derivatives}.
10\textsuperscript{a} edición, Pearson, Boston, 2017.

\bibitem{glasserman}
P. Glasserman.
\emph{Monte Carlo Methods in Financial Engineering}.
Springer, Nueva York, 2003.

\bibitem{broadie-glasserman-kou}
M. Broadie, P. Glasserman, S. Kou.
A continuity correction for discrete barrier options.
\emph{Mathematical Finance}, 7 (4), 325--349, 1997.

\bibitem{bouchaud-potters}
J.-P. Bouchaud, M. Potters.
\emph{Theory of Financial Risk and Derivative Pricing: From
Statistical Physics to Risk Management}.
2\textsuperscript{a} edición, Cambridge University Press,
Cambridge, 2003.

\bibitem{tfgm}
R. Cepeda Rocandio.
\emph{Valoración de opciones europeas: del modelo binomial al modelo
de Black--Scholes}.
Trabajo Fin de Grado en Matemáticas, Universidad de Oviedo, 2027.

\end{thebibliography}


\end{document}
